\section{Real structures, fixed groups, and the moduli comparison}
\label{sec:C}
Here $\Pi=\pi_1^{\mathrm{top}}(U(\mathbb C))$ is discrete,
$\widehat\Pi=\pi_1(U_{\mathbb C})$ is profinite, and
$c$ denotes the nonidentity element of $G_{\mathbb R}=C_2$.
The topological extension $E^{\mathrm{top}}$ consists of
lifts of the identity and complex conjugation to the universal
cover. Finite equivariant covering comparison identifies its
profinite completion with $E^{\mathrm{et}}=\pi_1(U_{\mathbb R})$.
Section classes are taken modulo $\Pi$ or $\widehat\Pi$,
respectively. The maps to automorphism-group quotients are
specified separately.

\subsection{The two classification clauses of Lead 3}

\begin{proposition}[C.1 --- Known, with reduction: the corrected surface count]
\label{C-classification}
For a compact connected oriented surface $X$ of genus $g$,
orientation-reversing involutions are classified, up to
orientation-preserving homeomorphism, by the quotient
surface $Y=X/\langle c\rangle$. The nonorientable quotients
have nonorientable genus $q\geq1$ and $r\geq0$ boundary
components, with $g=q+r-1$; the orientable quotients have
genus $h\geq0$ and $r\geq1$, with $g=2h+r-1$.
There are $g+2+\lfloor g/2\rfloor$ types in all.
The $g+1$ types counted in the original text of scan 439
are exactly the nonorientable-quotient types.
Its margin supplies the remaining types.
\end{proposition}
\begin{proof}
A collar of a fixed circle has reflection coordinates, so the
fixed circles give precisely the boundary components of the
quotient. Away from them the quotient map has degree two.
The oriented double of a nonorientable bordered surface,
and the double of an orientable bordered surface along its
nonempty boundary, recover $X$ with its involution.
These constructions are inverse on equivariant homeomorphism
classes; they type the morphism implicit on scan 437.
Euler characteristics give $2-2g=2(2-q-r)$ and
$2-2g=2(2-2h-r)$ respectively. The classification of compact
surfaces gives the stated ranges and counts.
These are also the restrictions of Gross--Harris,
Proposition 3.1, with the realization statement immediately
following it~\cite{GrossHarris1981}.
Grothendieck's $\gamma$ counts additional crosscaps after
the initial projective plane, so $q=\gamma+1$; his relation
$\gamma+r=g$ is the first formula. The margin on
folio 222, scan 439, is mathematically necessary.
\end{proof}

\begin{proposition}[C.2 --- Known and Refuted in the indicated senses: abelian subgroups and strips]
\label{C-abelian}
For the orientation-preserving mapping class group of a
surface of hyperbolic type $(g,n)$, every free abelian
subgroup has rank at most $3g-3+n$. A pants decomposition
gives that rank through its Dehn twists.
The assertion that all maximal free abelian subgroups
are generated by strips is false.
\end{proposition}
\begin{proof}
The bound is Birman--Lubotzky--McCarthy, Theorem A~\cite{BLM1983}.
The twists around the $3g-3+n$ pants curves commute.
Intersection arcs crossing one cuff and missing the others
show independence of their powers.
For the proposed identification on scan 470, use the
once-punctured torus and its homological identification
with $\mathrm{SL}_2(\mathbb Z)$. The matrix
\[
 A=\begin{pmatrix}2&1\\1&1\end{pmatrix}
\]
has characteristic polynomial $t^2-3t+1$. Its rational
commutant is $\mathbb Q[A]$, a real quadratic field.
The determinant-one integral units in this order have
rank one, so a maximal free abelian subgroup containing
$A$ has rank one. Every nontrivial Dehn twist has
unipotent homological action. No nonidentity unipotent
matrix commutes with $A$: its unique rational eigenline
would have to be one of the two irrational eigenlines
of $A$. Thus this subgroup contains no strip twist.
The characteristic polynomial and commutant equations
are checked in the exact program. This gives an
explicit homological comparison for the claim.
\end{proof}

\subsection{C(a): section classes and profinite centralizers}

\begin{theorem}[C.3 --- Known, with proof in his language: discrete real sections]
\label{C-discrete}
For every hyperbolic real curve $U$, the map
\[
 \pi_0(U(\mathbb R))\longrightarrow
 \operatorname{Sect}(E^{\mathrm{top}}\to C_2)/\Pi
\]
is bijective. For the lift $\sigma_I$ attached to a component $I$,
\[
 \Pi^{\sigma_I}=
 \begin{cases}
  1,& I\text{ is an interval},\\
  \langle a_I\rangle\simeq\mathbb Z,&I\text{ is an oval}.
 \end{cases}
\]
In the second case $a_I$ is the oval loop; choosing its
orientation accounts for $a_I\leftrightarrow a_I^{-1}$.
\end{theorem}
\begin{proof}
The hyperbolic universal cover is the disc.
An order-two lift of complex conjugation is an
orientation-reversing hyperbolic isometry, hence reflection
in one geodesic. Its fixed geodesic is connected and projects
onto one real component. The lift determined by a real
point fixes a point over it. All points in one component
therefore give the same conjugacy class, and one geodesic
cannot project into two components. This proves the bijection.
Deck transformations commuting with the reflection preserve
its fixed geodesic. Their restriction is the deck group of
the covering of that real component: trivial for an interval
and infinite cyclic for an oval. Restriction is injective
since a deck transformation fixing a point is the identity.
This proves the surjectivity in (14) as well as (6)--(10),
\S33, folios 223--225, scans 441--445. It is the
specialization of Mochizuki's Theorem 3.6 and
Corollary 3.7~\cite{MochizukiTopics}.
\end{proof}

\begin{theorem}[C.4 --- Proved: the profinite fixed group in every genus]
\label{C-profinite-centralizer}
For every hyperbolic real curve $U$, proper or affine,
completion gives a bijection of the classes in
Theorem~\ref{C-discrete} with
$\operatorname{Sect}(\pi_1(U)\to C_2)/\widehat\Pi$.
For every component $I$,
\[
 \widehat\Pi^{\sigma_I}=
 \begin{cases}
  1,&I\text{ is an interval},\\
  \overline{\langle a_I\rangle}\simeq\widehat{\mathbb Z},
       &I\text{ is an oval}.
 \end{cases}
\]
No bound on the genus of $U$ is required.
\end{theorem}
\begin{proof}
We verify the hypotheses of Boggi--Zalesskii, Theorem A,
for $E^{\mathrm{top}}$ and the prime two~\cite{BoggiZalesskii2024}.
Its index-two subgroup $\Pi$ is a free group of finite rank
or a closed orientable surface group. These groups are good.
For the required cohomological verification, first trivialize
a finite coefficient module on a finite cover. A degree-one
class is killed by the kernel of its finite-valued cocycle.
For a closed surface a degree-two class is killed on a cover
of degree divisible by the exponent of the coefficient
module: restriction on the fundamental class multiplies
by the degree. Higher cohomology vanishes. For a free
group only degree one occurs. The finite-cover criterion
for goodness applies. Comparing the two
Lyndon--Hochschild--Serre spectral sequences for the
index-two extension proves goodness of $E^{\mathrm{top}}$.
Finite coefficients have finite cohomology in every degree,
and the finite-index surface subgroup has cohomological
dimension at most two. Thus finite virtual
two-cohomological type holds.
Residual finiteness of the surface or free subgroup,
followed by the intersection of the finitely many
translates of a separating normal subgroup, proves
residual finiteness of the extension. For an arbitrary
discrete continuous $\mathbb F_2[\widehat E]$-module,
take the union of its finite submodules. The finite-type
resolutions make cohomology commute with this filtered
union, as required in the source's statement.

The cited theorem gives bijectivity for conjugacy classes
of order-two subgroups and identifies their completed
centralizers with closures of discrete centralizers.
The geometric group is torsion-free by finite
cohomological dimension. An $E$-conjugator between
two sections can be made geometric by multiplication
by the source involution. We therefore obtain exactly
the required conjugacy convention and
$C_{\widehat E}(\sigma_I)=\overline{C_E(\sigma_I)}$.
Intersect with the open subgroup $\widehat\Pi$ and use
Theorem~\ref{C-discrete}.

We check the full topology on an oval subgroup.
For each $m$, the essential simple loop $a_I$
has order divisible by $m$ in a finite quotient of $\Pi$.
For a nonseparating loop use homology modulo $m$.
For a separating loop cut along it. On a positive-genus
side assign one handle pair the permutations
$A=(1\,2\,\cdots\,m)$ and the interchange $B$ of
two blocks of size $m$. Then $[A,B]$ is an
$m$-cycle on the first block times the inverse cycle
on the second; it has order $m$. Use its inverse
on the other side. On a genus-zero side there is
another puncture; assign its loop that permutation
and the other loops the identity. These representations
glue by the boundary relation. An essential oval
has no unpunctured disc side. Thus every modulus
is detected, and its closure is $\widehat{\mathbb Z}$.
\end{proof}

\begin{proposition}[C.5 --- Refuted and repaired: the two signs at a real cusp]
\label{C-cusp-packet}
For a real cusp $d$ with decomposition extension
$1\to\widehat{\mathbb Z}(1)\to D_d\to C_2\to1$,
the local section-class set is a torsor under
$\mathbb Z/2$. The map into the global section set
sends its two classes to the components adjacent
to the two germs at $d$. These components can coincide.
The assertion of two distinct global classes for
every real cusp is false.
\end{proposition}
\begin{proof}
Conjugation acts by $-1$ on $\widehat{\mathbb Z}(1)$, so
\[
 H^1(C_2,\widehat{\mathbb Z}(1))
 =\widehat{\mathbb Z}/2\widehat{\mathbb Z}=\mathbb Z/2.
\]
The real Kummer map is the sign map. The two parameter
signs approach the two real germs; transport to a nearby
point identifies the tangential section with its component
class by Theorem~\ref{C-profinite-centralizer}.

Take $E:y^2=x^3+x$ and $U=E\setminus\{O\}$.
The cubic has one real root, so $E(\mathbb R)$ is one
oval. Deleting $O$ leaves one interval. The two
elements of $\operatorname{Sect}(D_O\to C_2)/
\widehat{\mathbb Z}(1)$ therefore have the same
image in $\operatorname{Sect}(\pi_1(U)\to C_2)/
\pi_1(U_{\mathbb C})$. This is the exact map and
fibre giving the counterexample.
\end{proof}

\begin{proposition}[C.6 --- Proved: the incidence graph recovers the real boundary order]
\label{C-boundary-order}
The arithmetic extension of a hyperbolic real curve,
together with its peripheral conjugacy classes,
recovers the real cusps, the component of $X(\mathbb R)$
containing each, and their cyclic order up to reversal.
This includes an oval carrying just one real cusp.
\end{proposition}
\begin{proof}
A cusp is real exactly when its inertia conjugacy class
is fixed by $C_2$; other cusps form two-element orbits.
For each real cusp take its two-element local section
torsor and its map to the global section set.
Make an unoriented edge for the cusp, whose ends are
the two image classes. Interval vertices are also
recognized by their trivial fixed group in
Theorem~\ref{C-profinite-centralizer}.
Every such interval has two ends, so the resulting
graph is a disjoint union of cycles. A one-edge loop
occurs precisely when the two local signs have the
same global image. Traversing a cycle lists the
cusps of one compact oval in their cyclic order,
with reversal as its ambiguity. Vertices with
fixed group $\widehat{\mathbb Z}$ are the ovals
with no real cusp. Nonreal pairs give no graph ends.
This answers the margin of \S33, folio 221, scan 437.
\end{proof}

\subsection{C(b)--(c): arithmetic conjugations and the tripod}

\begin{proposition}[C.7 --- Proved: prediscretizations and archimedean places]
\label{C-prediscretizations}
For the arithmetic prediscretization torsor $P_a$
of scans 377--378, (29)--(30) are the maps
\[
 e\longmapsto e^{-1}ce,\qquad
 P_a/\{1,c\}\longrightarrow
 \{\text{complex conjugations in }\Gamma_a\}.
\]
The second is a bijection. Under the faithful
identification $\Gamma_a\simeq G_{\mathbb Q}$, its
image is every element of order two and it
identifies $P_a/\{\pm1\}$ with the archimedean
places of $\overline{\mathbb Q}$.
The cyclotomic-kernel orbits on this set are
$\widehat{\mathbb Z}^{\,\times}/\{\pm1\}$.
\end{proposition}
\begin{proof}
The quotient construction of LM3 A.8 and B.6 makes
the oriented arithmetic prediscretizations a torsor
under the arithmetic image. Choosing a marking
identifies it with embeddings of the fixed algebraic
closure into $\mathbb C$. Pulling back conjugation
gives the displayed involution. Reversing orientation
composes the embedding with $c$ and has the same
image. Before faithfulness this gives exactly the
conjugates of $c$ in that image.

We compute the stabilizer. An automorphism of
$\overline{\mathbb Q}$ commuting with $c$ preserves
its fixed field $R_{\rm alg}$ of real algebraic
numbers. Positivity is the set of nonzero squares,
so it preserves the order. Since it fixes
$\mathbb Q$ and each real algebraic number is
determined by its rational cut, it is the
identity on $R_{\rm alg}$. As
$\overline{\mathbb Q}=R_{\rm alg}(i)$, it is
$1$ or $c$. Thus $C_{G_{\mathbb Q}}(c)=\{1,c\}$.
An arbitrary involution has a real closed fixed
field by Artin--Schreier. Its ordering gives an
embedding into $\mathbb R$; extension to the
algebraic closure conjugates it to $c$.
Embeddings modulo conjugation are the archimedean
places. Hoshi--Mochizuki, Corollary 6.4, supplies
the faithful arithmetic image used here; this
quotient argument needs no assertion about the
geometric congruence kernel.
Finally $\chi$ is onto and $\chi(c)=-1$.
The double-coset calculation gives the last formula.
\end{proof}

\begin{proposition}[C.8 --- Proved: restriction and the six tripod germs]
\label{C-tripod}
Restriction of a number-field section to a real
decomposition group is its real component section.
For $B=\mathbb P^1_{\mathbb Q}\setminus\{0,1,\infty\}$,
choose the unit tangential base point at $0$ toward $1$.
There are free profinite generators $x,y$ with
$c(x)=x^{-1}$ and $c(y)=y^{-1}$, and the three
real section involutions are
\[
 \tau_{(0,1)}=c,\qquad
 \tau_{(-\infty,0)}=xc,\qquad
 \tau_{(1,\infty)}=y^{-1}c .
\]
The cusp germs have the following global images:
\[
\begin{array}{c|c|c}
\text{cusp}&\text{local parameter sign}&\text{involution}\\ \hline
0&t>0&c\\
0&t<0&xc\\
1&t-1<0&c\\
1&t-1>0&y^{-1}c\\
\infty&1/t>0&y^{-1}c\\
\infty&1/t<0&xc
\end{array}
\]
These are actual equalities after the stated
path choices; forgetting them gives geometric
conjugacy classes.
\end{proposition}
\begin{proof}
The base-change square of arithmetic fundamental
groups over $K$ and $K_v=\mathbb R$ is a pullback
over $G_K$ and $G_{\mathbb R}$. Pulling back the
section and applying
Theorem~\ref{C-profinite-centralizer} proves compatibility.
For the formulas use the real half-plane presentation
\[
 E^{\mathrm{top}}
 =\langle a,b,d\mid a^2=b^2=d^2=1\rangle,\quad
 c=a,\quad x=ba,\quad y=ad .
\]
The geometric subgroup is the kernel sending all
three reflections to $c$. The reflected sides
project to the three intervals. The third
peripheral generator is $z=(xy)^{-1}=db$.
The reflections at infinity are $d$ and
$z^{-1}d=b$; at the other cusps they are
$a,b$ and $a,d$. This proves the six entries.
The program checks the free-product identities.
The completion bijection shows the three
classes remain distinct profinitely.
\end{proof}

\begin{proposition}[C.9 --- Proved: an explicit class-two real obstruction]
\label{C-tripod-class-two}
The abelianized real tripod has four section classes
in $(\mathbb Z/2)^2$. Exactly
$(0,0),(1,0),(0,1)$ lift fully. The fourth already
fails in a finite nilpotent quotient of class two.
\end{proposition}
\begin{proof}
Inversion gives $H^1(C_2,\widehat{\mathbb Z}(1)^2)
=(\mathbb Z/2)^2$, and Proposition~\ref{C-tripod}
gives the first three classes. For the fourth use
\[
 F_2\longrightarrow
 D_8=\langle r,t\mid r^4=t^2=1,\ trt=r^{-1}\rangle,
 \qquad x\mapsto t,\quad y\mapsto rt .
\]
The real action becomes trivial since the images
of $x,y$ are involutions. Thus the extension is
$D_8\times C_2$. The geometric coordinates above
the fourth abelian class are $r$ and $r^3$;
both square to $r^2\ne1$, preventing a section.
The program exhaustively checks that coset.
The quotient maps to this class-two extension
and the abelian extension specify the comparison.
\end{proof}

\subsection{C(d): the meaning of the component counts}

\begin{proposition}[C.10 --- Proved: effective lifts and the centralizer index]
\label{C-moduli-count}
Fix an orientation-reversing mapping class $\sigma$
of a compact genus-$g$ surface, $g\geq2$, with
$r$ real ovals. The fibre over $\sigma$ of
order-two lifts in $\operatorname{Aut}(\Pi_g)$,
modulo inner automorphisms, has $r$ elements.
The centralizer of $\sigma$ in the
orientation-preserving outer mapping class
group acts transitively on this fibre.
For one lift $u$ its stabilizer fits into
\[
\begin{split}
 1\longrightarrow\Pi_g^u
 &\longrightarrow C_{\operatorname{Aut}^{+}(\Pi_g)}(u)\\
 &\longrightarrow
 \operatorname{Stab}_{C_{\operatorname{Out}^{+}(\Pi_g)}(\sigma)}
       ([u])\longrightarrow1.
\end{split}
\]
The last group has index $r$ in the outer centralizer
and $\Pi_g^u\simeq\mathbb Z$.
Thus the count $r$ on scans 445--448 is the size of
this fibre before the centralizer acts.
The unmarked real moduli stack has
$g+2+\lfloor g/2\rfloor$ components; with one
labelled real marked point it has
$g+1+\lfloor g/2\rfloor$.
\end{proposition}
\begin{proof}
Apply Theorem~\ref{C-discrete} to a representative
of $\sigma$. Its ovals give exactly the effective
lifts and their fixed groups. Homeomorphisms of the
quotient surface permute its boundary circles
transitively. Their orientation-preserving lifts
to the double centralize $\sigma$, proving
transitivity on the lift classes.
An outer centralizer element stabilizes $[u]$
precisely when a representative can be modified
by an inner automorphism to commute with $u$.
The kernel consists of inner automorphisms by
$\Pi_g^u$. This proves the sequence and the
index. The index assertion on scan 447 must
therefore be applied to the displayed image,
with its infinite cyclic kernel retained.

Mochizuki's Corollary 3.8 identifies real moduli
components with these orientation-reversing
mapping-class conjugacy classes.
Proposition~\ref{C-classification} gives the
unmarked count. One real marked point requires
$r>0$, and transitivity gives one component
for each such type. Exactly the fixed-point-free
type is removed. This quotient of the $r$ lift
classes to one centralizer orbit is the
precise comparison needed for scans 445--448.
\end{proof}

\begin{proposition}[C.11 --- Proved: a finite formula for labelled real markings]
\label{C-marked-count}
For $g\geq2$, $n>0$, the components of the real
stack with $n$ individually labelled real marked
points are classified by an allowed type $(r,a)$
of Proposition~\ref{C-classification}, $r>0$,
and the following data.
Partition $\{1,\dots,n\}$ into at most $r$
nonempty blocks and put a cyclic order on each.
For a nonorientable quotient ($a=1$) identify
each cyclic order with its reverse independently;
for an orientable quotient ($a=0$) identify
the simultaneous reversal of all orders.
For a partition of block sizes $m_1,\dots,m_k$
the contributions are
\[
 \prod_{j=1}^k b(m_j)\quad(a=1),\qquad
 \frac{\prod_{j=1}^k(m_j-1)!+
              {\bf1}_{\max m_j\leq2}}2\quad(a=0),
\]
where $b(1)=b(2)=1$ and $b(m)=(m-1)!/2$ for $m\geq3$.
Sum over the allowed types and partitions.
\end{proposition}
\begin{proof}
Marked real points lie on boundary circles of
the quotient. Boundary permutations account
for the blocks, with cyclic orders within
them. All boundary permutations are realized.
On an orientable quotient a homeomorphism
preserves all induced boundary orientations,
or reverses all of them. On a nonorientable
quotient, moving a small boundary disc along
a one-sided loop and back reverses that
disc's orientation and fixes the other
boundary circles. These moves give precisely
the independent reversals.

Circle configuration spaces with prescribed
cyclic order are connected, as are the real
Teichmüller loci for a fixed involution.
Thus these homeomorphism orbits give exactly
the connected components.
A labelled cyclic order on $m$ points has
$(m-1)!$ possibilities; reversal fixes one
precisely for $m\leq2$. Burnside's formula
gives both counts.
For the manuscript's unordered punctures,
the complete topological data instead
allow a finite subset in the boundary
and interior of the quotient; interior
points lift to conjugate pairs. This is
the double construction of scan 437.
The finite permutation quotient of the
labelled moduli stack is the comparison
map. The numerical formula here specifies
its individually real labelled objects.
\end{proof}

\begin{proposition}[C.12 --- Open: full profinite moduli in genus at least three]
\label{C-high-genus-open}
The classification of all involutions over complex
conjugation in the full profinite extended mapping
class group, with complete centralizers, is not
established here for $g\geq3$. The smallest
uncompleted type in this calculation is $(3,0)$.
The exact fibre calculation below records the
obstruction reached by the attempted routes.
\end{proposition}
\begin{proof}[Established part and attempted routes]
Mochizuki's Corollary 3.15 proves the full
statement for $g\leq2$. The goodness proof
in Theorem~\ref{C-profinite-centralizer}
uses a free or surface subgroup of finite
index. The group here is connected to
a surface group by the point-pushing
map in the Birman exact sequence.
Induction along this sequence changes
the number of markings and retains
the genus, so it does not give the
missing genus-three cohomology
comparison.

We checked Holzschuh's Theorems A and
B~\cite{Holzschuh2025}. Theorem A
concerns the geometrically pro-two
étale homotopy type of an equivariantly
triangulable scheme; Theorem B assumes
geometric étale nilpotence for the
full homotopy type. The comparison
relevant to a level cover $X$ is
\[
 \operatorname{Et}(X)\longrightarrow
 B\pi_1^{\mathrm{et}}(X).
\]
Lifting a section on the right
requires the successive Postnikov
obstructions to vanish. Their
targets are
$H^{q+1}(C_2,\pi_q(\operatorname{Et}(X_{\mathbb C})))$,
with the action furnished by the
partial lift, for $q\geq2$.
The source does not compute these
obstructions for a level cover
of $\mathcal M_3$. The need to
retain higher homotopy is visible
already from the unmarked group:
the genus-at-least-three mapping
class group is perfect, so its
maximal pro-two quotient is trivial.

For the congruence route write
\[
 1\to K_{g,n}\to\widehat\Gamma^{\,\pm}_{g,n}
       \to\check\Gamma^{\,\pm}_{g,n}\to1 .
\]
Fix a topological involution $\tau$.
The profinite lifts of its congruence
image, modulo $K_{g,n}$-conjugacy,
form exactly $H^1(C_2,K_{g,n})$
for conjugation by $\tau$.
Indeed $k\tau$ squares to one
exactly when $k\,\tau(k)=1$,
and conjugation gives the
nonabelian coboundary relation.
The kernel on centralizers is
$K_{g,n}^{\tau}$.
We have not established the required
orbit calculation or the fixed
kernel for $(g,n)=(3,0)$.

The symmetric congruence theorem
of Boggi~\cite{Boggi2026} applies
to centralizers of finite
orientation-preserving actions
with quotient genus at most two.
We checked its actual map
$\widehat{\Gamma(S)^G}\to
\operatorname{Out}(\widehat\Pi)$.
Here the reflection lies in
the orientation extension, and
the required centralizer is
inside the full completed group.
The displayed $H^1$ and fixed
kernel remain to be determined.
A next concrete step is the
two-cohomology comparison for
a genus-three level subgroup,
or a direct calculation of
these two kernel invariants.
This record asserts no
counterexample to the
genus-three statement.
\end{proof}


